% Lösungen Zone (zusammenbau v0.1)
\einheitenkopf[][Kennst du schon]{Das kennst du schon}

Zuordnung: 1 → die Prüfregeln aller Einheiten · 2 → das Einsetzen in Einheit 3 und 4 · 3 → die Merkregel in Einheit 3 und die Kreisbahn-Begründung · 4 → Einheit 2 und 4 · 5 → die Scharen in Einheit 3 · 6 → die Bereichsprüfungen in Einheit 4 · 7 → die Prüfregeln aller Einheiten · 8 → die Prüfregeln aller Einheiten

\erg{1}{a) $\vec{n} = (2 | 3 | -1)$ \quad b) $\vec{n} = (4 | -0,5 | 3)$}
\erg{2}{a) $(3 | 4 | 2)$ \quad b) $(1 + 2t | -2 + t | 3 - t)$}
\erg{3}{a) $7$ \quad b) $-6$}
\erg{4}{a) $k = -2$ \quad b) $t = 0,5$}
\erg{5}{a) $a + 2$ \quad b) $t = \frac{4}{a + 2}$}
\erg{6}{a) ja – jede Koordinate liegt in ihren Schranken. \quad b) nein – die z-Koordinate ist negativ, der Punkt liegt unter dem Boden.}
\erg{7}{a) Tim hat die fehlende Variable weggelassen; im Raum hat der Normalenvektor drei Komponenten: $\vec{n} = (2 | -1 | 0)$.}
\erg{8}{a) $\vec{n} = (5 | 0 | 1)$}
